%%% The main file. It contains definitions of basic parameters and includes all other parts.

%% Settings for single-side (simplex) printing
% Margins: left 40mm, right 25mm, top and bottom 25mm
% (but beware, LaTeX adds 1in implicitly)
\documentclass[12pt,a4paper]{report}
\setlength\textwidth{145mm}
\setlength\textheight{247mm}
\setlength\oddsidemargin{15mm}
\setlength\evensidemargin{15mm}
\setlength\topmargin{0mm}
\setlength\headsep{0mm}
\setlength\headheight{0mm}
% \openright makes the following text appear on a right-hand page
\let\openright=\clearpage

%% Settings for two-sided (duplex) printing
% \documentclass[12pt,a4paper,twoside,openright]{report}
% \setlength\textwidth{145mm}
% \setlength\textheight{247mm}
% \setlength\oddsidemargin{14.2mm}
% \setlength\evensidemargin{0mm}
% \setlength\topmargin{0mm}
% \setlength\headsep{0mm}
% \setlength\headheight{0mm}
% \let\openright=\cleardoublepage

%% Generate PDF/A-2u
\let\pdfcreationdate=\creationdate
\usepackage[a-2u]{pdfx}

%% Character encoding: usually latin2, cp1250 or utf8:
\usepackage[utf8]{inputenc}

%% Prefer Latin Modern fonts
\usepackage{lmodern}

%% Further useful packages (included in most LaTeX distributions)
\usepackage{amsmath}        % extensions for typesetting of math
\usepackage{amsfonts}       % math fonts
\usepackage{amsthm}         % theorems, definitions, etc.
\usepackage{bbding}         % various symbols (squares, asterisks, scissors, ...)
\usepackage{bm}             % boldface symbols (\bm)
\usepackage{graphicx}       % embedding of pictures
\usepackage{fancyvrb}       % improved verbatim environment
\usepackage{natbib}         % citation style AUTHOR (YEAR), or AUTHOR [NUMBER]
\usepackage[nottoc]{tocbibind} % makes sure that bibliography and the lists
			    % of figures/tables are included in the table
			    % of contents
\usepackage{dcolumn}        % improved alignment of table columns
\usepackage{booktabs}       % improved horizontal lines in tables
\usepackage{paralist}       % improved enumerate and itemize
\usepackage{xcolor}         % typesetting in color

%%% Basic information on the thesis

% Thesis title in English (exactly as in the formal assignment)
\def\ThesisTitle{Command-line tool lsql-csv for CSV files processing}

% Author of the thesis
\def\ThesisAuthor{Přemysl Šťastný}

% Year when the thesis is submitted
\def\YearSubmitted{2024}

% Name of the department or institute, where the work was officially assigned
% (according to the Organizational Structure of MFF UK in English,
% or a full name of a department outside MFF)
\def\Department{Department of Applied Mathematics}

% Is it a department (katedra), or an institute (ústav)?
\def\DeptType{Department}

% Thesis supervisor: name, surname and titles
\def\Supervisor{doc. Mgr. Jan Hubička, Ph.D.}

% Supervisor's department (again according to Organizational structure of MFF)
\def\SupervisorsDepartment{Department of Applied Mathematics}

% Study programme and specialization
\def\StudyProgramme{Computer Science}
\def\StudyBranch{General Computer Science}

% An optional dedication: you can thank whomever you wish (your supervisor,
% consultant, a person who lent the software, etc.)
\def\Dedication{%
Dedication.
}

% Abstract (recommended length around 80-200 words; this is not a copy of your thesis assignment!)
\def\Abstract{%
  lsql-csv is a tool for small CSV file data querying from a shell with short
  queries. It makes it possible to work with small CSV files like with a read-only
  relational databases. The tool implements a new language LSQL similar to SQL,
  specifically designed for working with CSV files in shell. 
  LSQL aims to be a more lapidary language than SQL.
  Its design purpose is to enable its user to quickly write simple queries directly
  to the terminal.
}

% 3 to 5 keywords (recommended), each enclosed in curly braces
\def\Keywords{%
  {relational database}, {csv}, {sql}, {Haskell}, {query language}, {unix philosophy}
}

%% The hyperref package for clickable links in PDF and also for storing
%% metadata to PDF (including the table of contents).
%% Most settings are pre-set by the pdfx package.
\hypersetup{unicode}
\hypersetup{breaklinks=true}
\hypersetup{hidelinks}

% Definitions of macros (see description inside)
%%% This file contains definitions of various useful macros and environments %%%
%%% Please add more macros here instead of cluttering other files with them. %%%

%%% Minor tweaks of style

% These macros employ a little dirty trick to convince LaTeX to typeset
% chapter headings sanely, without lots of empty space above them.
% Feel free to ignore.
\makeatletter
\def\@makechapterhead#1{
  {\parindent \z@ \raggedright \normalfont
   \Huge\bfseries \thechapter. #1
   \par\nobreak
   \vskip 20\p@
}}
\def\@makeschapterhead#1{
  {\parindent \z@ \raggedright \normalfont
   \Huge\bfseries #1
   \par\nobreak
   \vskip 20\p@
}}
\makeatother

% This macro defines a chapter, which is not numbered, but is included
% in the table of contents.
\def\chapwithtoc#1{
\chapter*{#1}
\addcontentsline{toc}{chapter}{#1}
}

% Draw black "slugs" whenever a line overflows, so that we can spot it easily.
\overfullrule=1mm

%%% Macros for definitions, theorems, claims, examples, ... (requires amsthm package)

\theoremstyle{plain}
\newtheorem{thm}{Theorem}
\newtheorem{lemma}[thm]{Lemma}
\newtheorem{claim}[thm]{Claim}

\theoremstyle{plain}
\newtheorem{defn}{Definition}

\theoremstyle{remark}
\newtheorem*{cor}{Corollary}
\newtheorem*{rem}{Remark}
\newtheorem*{example}{Example}

%%% An environment for proofs

\newenvironment{myproof}{
  \par\medskip\noindent
  \textit{Proof}.
}{
\newline
\rightline{$\qedsymbol$}
}

%%% An environment for typesetting of program code and input/output
%%% of programs. (Requires the fancyvrb package -- fancy verbatim.)

\DefineVerbatimEnvironment{code}{Verbatim}{fontsize=\small, frame=single}

%%% The field of all real and natural numbers
\newcommand{\R}{\mathbb{R}}
\newcommand{\N}{\mathbb{N}}

%%% Useful operators for statistics and probability
\DeclareMathOperator{\pr}{\textsf{P}}
\DeclareMathOperator{\E}{\textsf{E}\,}
\DeclareMathOperator{\var}{\textrm{var}}
\DeclareMathOperator{\sd}{\textrm{sd}}

%%% Transposition of a vector/matrix
\newcommand{\T}[1]{#1^\top}

%%% Various math goodies
\newcommand{\goto}{\rightarrow}
\newcommand{\gotop}{\stackrel{P}{\longrightarrow}}
\newcommand{\maon}[1]{o(n^{#1})}
\newcommand{\abs}[1]{\left|{#1}\right|}
\newcommand{\dint}{\int_0^\tau\!\!\int_0^\tau}
\newcommand{\isqr}[1]{\frac{1}{\sqrt{#1}}}

%%% Various table goodies
\newcommand{\pulrad}[1]{\raisebox{1.5ex}[0pt]{#1}}
\newcommand{\mc}[1]{\multicolumn{1}{c}{#1}}

%%% Various code formatting goodies
\def\icode#1{\texttt{#1}}


% Title page and various mandatory informational pages
\begin{document}
%%% Title page of the thesis and other mandatory pages

%%% Title page of the thesis

\pagestyle{empty}
\hypersetup{pageanchor=false}
\begin{center}

\centerline{\mbox{\includegraphics[width=166mm]{../img/logo-en.pdf}}}

\vspace{-8mm}
\vfill

{\bf\Large BACHELOR THESIS}

\vfill

{\LARGE\ThesisAuthor}

\vspace{15mm}

{\LARGE\bfseries\ThesisTitle}

\vfill

\Department

\vfill

{
\centerline{\vbox{\halign{\hbox to 0.45\hsize{\hfil #}&\hskip 0.5em\parbox[t]{0.45\hsize}{\raggedright #}\cr
Supervisor of the bachelor thesis:&\Supervisor \cr
\noalign{\vspace{2mm}}
Study programme:&\StudyProgramme \cr
\noalign{\vspace{2mm}}
Study branch:&\StudyBranch \cr
}}}}

\vfill

% Zde doplňte rok
Prague \YearSubmitted

\end{center}

\newpage

%%% Here should be a bound sheet included -- a signed copy of the "bachelor
%%% thesis assignment". This assignment is NOT a part of the electronic
%%% version of the thesis. DO NOT SCAN.

%%% A page with a solemn declaration to the bachelor thesis

\openright
\hypersetup{pageanchor=true}
\pagestyle{plain}
\pagenumbering{roman}
\vglue 0pt plus 1fill

\noindent
I declare that I carried out this bachelor thesis independently, and only with the cited
sources, literature and other professional sources. It has not been used to obtain another
or the same degree.

\medskip\noindent
I understand that my work relates to the rights and obligations under the Act No.~121/2000 Sb.,
the Copyright Act, as amended, in particular the fact that the Charles
University has the right to conclude a license agreement on the use of this
work as a school work pursuant to Section 60 subsection 1 of the Copyright~Act.

\vspace{10mm}

\hbox{\hbox to 0.5\hsize{%
In \hbox to 6em{\dotfill} date \hbox to 6em{\dotfill}
\hss}\hbox to 0.5\hsize{\dotfill\quad}}
\smallskip
\hbox{\hbox to 0.5\hsize{}\hbox to 0.5\hsize{\hfil Author's signature\hfil}}

\vspace{20mm}
\newpage

%%% Dedication

\openright

\noindent
\Dedication

\newpage

\openright

\noindent
\Thanks

\newpage

%%% Thanks

%%% Mandatory information page of the thesis

\openright

\vbox to 0.5\vsize{
\setlength\parindent{0mm}
\setlength\parskip{5mm}

Title:
\ThesisTitle

Author:
\ThesisAuthor

\DeptType:
\Department

Supervisor:
\Supervisor, \SupervisorsDepartment

Abstract:
\Abstract

Keywords:
\Keywords

\vss}

\newpage

\openright
\pagestyle{plain}
\pagenumbering{arabic}
\setcounter{page}{1}


%%% A page with automatically generated table of contents of the bachelor thesis

\tableofcontents

%%% Each chapter is kept in a separate file
\chapter*{Introduction}
\addcontentsline{toc}{chapter}{Introduction}

Database refers to a set of related data accessed through the use of a database management system \cite{enwiki-database}.
CSV files (Comma Separated Value files) are common way of exchanging and converting data between various spreadsheet programs \cite{rfc4180}.
Through this definition, we can see even a simple collection of CSV files accessed through some programs may be seen as an database itself.


SQL (Structured Query Language) is a language used to manage data, especially in a relational database management system \cite{enwiki-sql}.
It was first introduced at 1970s \cite{enwiki-sql} and is one of the most used query languages. 
Despite being standardized at 1987 by International Organization for Standardization \cite{ISO9075-1987}, 
there are virtually no implementations adhere to it fully \cite{enwiki-sql}.
Standard SQL is a typed language (Every data value belongs to some data type) \cite{ISO9075-2023} and the language
design is therefore not much suitable for type-less databases like a collection of CSV files.
Despite that, there are some implementations of SQL (eg. q \cite{q}), which tries to implement SQL on CSV files.

SQL itself requires large amount of text to be written for running even a simple queries and Unix ecosystem misses a tool\footnote{Or author don't know about it}, 
which would allowed running a short enough queries over CSV files with similar semantics like SQL has. And this is the reason, why lsql-csv was created.

Haskell is a language with great features for working with text \cite{practical-haskell} and therefore it was selected for the task of implementation of lsql-csv.

lsql-csv is a tool for CSV files data querying from shell with short queries. It makes possible to work with small CSV files like with read-only relational database.
The tool implements a new language LSQL similar to SQL, which is type-less, specifically designed for working with CSV files in shell.

\chapter{User documentation}
\textit{lsql-csv} is a tool for CSV file data querying from the shell with short queries. It makes it possible to work with small CSV files like with a read-only relational database.

The tool implements a new language LSQL similar to SQL, specifically designed for working with CSV files in shell. 

\section{Installation}
It is necessary, you had \textit{GHC} ($\geq 8 <9.29$) and Haskell packages \textit{Parsec} ($\geq 3.1$ $<3.2$), \textit{Glob} ($\geq 0.10$ $<0.11$), 
\textit{base} ($\geq 4.9$ $<4.17$), \textit{text} ($\geq 1.2$ $<1.3$) and \textit{containers} ($\geq 0.5$ $<0.7$)
 installed. (The package boundaries given are identical to cabal boundaries.) Run then:

\begin{verbatim}
make
sudo make install
\end{verbatim}

Now the lsql-csv is installed in \textit{/usr/local/bin}. If you want, you can specify \textit{INSTALL\_DIR} like:
\begin{verbatim}
sudo make INSTALL_DIR=/custom/install-folder install
\end{verbatim}
This will install the package into \textit{INSTALL\_DIR}.

If you have installed \textit{cabal}, you can alternatively run:
\begin{verbatim}
cabal install
\end{verbatim}
It will also install the dependencies for you.

\subsection{Running the unit tests}
If you want to verify, that the package has been compiled correctly, it is possible to test it by running:
\begin{verbatim}
make test
\end{verbatim}
This will run all unit tests for you.

\section{lsql-csv---quick introduction}
LSQL, the language of \textit{lsql-csv}, aims to be a more lapidary language than SQL. Its design purpose is to enable its user to quickly write simple queries directly to the terminal---its design purpose is therefore different from SQL, where the readability of queries is more taken into account than in LSQL.

\subsection{Examples}
One way to learn a new programming language is by understanding concrete examples of its usage. The following examples are written explicitly for the purpose of teaching a reader, how to use the tool \textit{lsql-csv} by showing him many examples of its usage.

The following examples might not be enough for readers, who don't know Unix/Linux scripting enough. If this is the case, please consider learning Unix/\allowbreak Linux scripting first before LSQL.

It is also advantageous to know SQL.

The following examples will be mainly about parsing of \textit{/etc/passwd} and parsing of \textit{/etc/group}. To make example reading more comfortable, we have added \textit{/etc/passwd} and \textit{/etc/group} column descriptions from man pages to the text.

File \textit{/etc/passwd} has the following columns\cite{passwd}:
\begin{enumerate}
    \item login name;
    \item optional encrypted password;
    \item numerical user ID;
    \item numerical group ID;    
    \item user name or comment field;
    \item user home directory;
    \item optional user command interpreter.
\end{enumerate}
File \textit{/etc/group} has the following columns\cite{group}:
\begin{enumerate}
    \item group name;
    \item password;
    \item numerical group ID;
    \item user list.
\end{enumerate}


\subsubsection{Hello World}
\begin{verbatim}
lsql-csv '-, &1.2 &1.1'
\end{verbatim}
This will print the second (\textit{\&1.2}) and the first column (\textit{\&1.1}) of csv file on stdin. If you know SQL, you can read it like \textit{from stdio S select S.second, S.first}.

Commands are split by commas into blocks. The first block is (and always is) the from block. There are file names or \textit{-} (stdin) separated by space. The second block is the select block, also separated by space.

For example:
\begin{verbatim}
lsql-csv '-, &1.2 &1.1' <<- EOF
World,Hello
EOF
\end{verbatim}
It returns:
\begin{verbatim}
Hello,World
\end{verbatim}


\subsubsection{Simple filtering}

\begin{verbatim}
lsql-csv -d: '-, &1.*, if &1.3>=1000' < /etc/passwd
\end{verbatim}
This will print lines of users whose UID $\geq 1000$. It can also be written as:
\begin{verbatim}
lsql-csv -d: 'p=/etc/passwd, p.*, if p.3 >= 1000'
    
lsql-csv -d: 'p=/etc/passwd, &1.*, if &1.3 >= 1000'

lsql-csv -d: '/etc/passwd, &1.*, if &1.3 >= 1000'
\end{verbatim}
The \textit{-d:} optional argument means the primary delimiter is \textit{:}. In previous examples we used overnaming, which allows us to give a data source file \textit{/etc/passwd} a name \textit{p}.

If you know SQL, you can read it as \textit{from /etc/passwd P select * where P.UID \textgreater= 1000}. As you can see, lsql style is more compressed than standard SQL.

The output might be:
\begin{verbatim}
nobody:x:65534:65534:nobody:/var/empty:/bin/false
me:x:1000:1000::/home/me:/bin/bash
\end{verbatim}

If you specify delimiter specifically for \textit{/etc/passwd}, the output will be a comma delimited.
\begin{verbatim}
lsql-csv '/etc/passwd -d:, &1.*, if &1.3 >= 1000'
\end{verbatim}
It might return:
\begin{verbatim}
nobody,x,65534,65534,nobody,/var/empty,/bin/false
me,x,1000,1000,,/home/me,/bin/bash
\end{verbatim}


\subsubsection{Named columns}
Let's suppose we have a file people.csv:
\begin{verbatim}
name,age                                                                                                                                                                                                                                   
Adam,21                                                                                                                                                                                                                                    
Petra,23                                                                                                                                                                                                                                   
Karel,25
\end{verbatim}

Now, let's get all the names of people in people.csv using the \textit{-n} named switch:
\begin{verbatim}
lsql-csv -n 'people.csv, &1.name'
\end{verbatim}
The output will be:
\begin{verbatim}
Adam
Petra
Karel
\end{verbatim}

As you can see, we can reference named columns by a name. Named switch \textit{-n} enables first-line headers.
If named columns are enabled, each column has two names under \textit{\&X} - the number name \textit{\&X.Y} and the actual name \textit{\&X.NAME}.

Now, we can select all columns with wildcard \textit{\&1.*}:
\begin{verbatim}
lsql-csv -n 'people.csv, &1.*'
\end{verbatim}
As the output, we get
\begin{verbatim}
Adam,21,21,Adam
Petra,23,23,Petra
Karel,25,25,Karel
\end{verbatim}

The output contains each column twice because wildcard \textit{\&1.*} was evaluated to \textit{\&1.1, \&1.2, \&1.age, \&1.name}.
How to fix it?
\begin{verbatim}
lsql-csv -n 'people.csv, &1.[1-9]*'
\end{verbatim}
The output is now:
\begin{verbatim}
Adam,21
Petra,23
Karel,25
\end{verbatim}

The command can also be written as
\begin{verbatim}
lsql-csv -n 'people.csv, &1.{1,2}'
lsql-csv -n 'people.csv, &1.{1..2}'
\end{verbatim}
The output will be in both cases still the same.


\subsubsection{Simple join}
Let's say, I am interested in the default group names of users. We need to join to tables: \textit{/etc/passwd} and \textit{/etc/group}. Let's do it.
\begin{verbatim}
lsql-csv -d: '/etc/{passwd,group}, &1.1 &2.1, if &1.4 == &2.3'
\end{verbatim}
What does \textit{/etc/\{passwd,group\}} mean? Basically, there are three types of expressions. Select, from and arithmetic expression. In all select and from expressions, you can use expansion and wildcards just like in bash\cite{bash-reference-manual}.

Finally, the output can be something like this:
\begin{verbatim}
root:root
bin:bin
daemon:daemon
me:me
\end{verbatim}
The first column is the name of a user and the second column is the name of its default group.


\subsubsection{Basic grouping}
Let's say, I want to count users using the same shell.
\begin{verbatim}
lsql-csv -d: 'p=/etc/passwd, p.7 count(p.3), by p.7'
\end{verbatim}
And the output?
\begin{verbatim}
/bin/bash:7
/bin/false:7
/bin/sh:1
/bin/sync:1
/sbin/halt:1
/sbin/nologin:46
/sbin/shutdown:1
\end{verbatim}

You can see here the first usage of \textit{by} block, which is equivalent of \textit{group by} in SQL.

\subsubsection{Basic sorting}
Let's say, you want to sort your users with UID greater than or equal to 1000 ascendingly.
\begin{verbatim}
lsql-csv -d: '/etc/passwd, &1.*, if &1.3 >= 1000, sort &1.3'
\end{verbatim}
The output might look like:
\begin{verbatim}
me1:x:1000:1000::/home/me1:/bin/bash
me2:x:1001:1001::/home/me2:/bin/bash
me3:x:1002:1002::/home/me3:/bin/bash
nobody:x:65534:65534:nobody:/var/empty:/bin/false
\end{verbatim}

The sort block is the equivalent of \textit{order by} in SQL.

If we wanted descendingly sorted output, we might create a pipe to the \textit{tac} command - the \textit{tac} command prints the lines in reverse order:
\begin{verbatim}
lsql-csv -d: '/etc/passwd, &1.*, if &1.3 >= 1000, sort &1.3' | tac
\end{verbatim}

\subsubsection{About nice outputs}
There is a trick, how to concatenate two values in a select expression: Write them without space.

But how does the interpreter know the ends of the value name or value expression? You must use quotes for it - quotes themselves can't be part of the value name.
As an example, let's try to format our basic grouping example.

Let's try it!
\begin{verbatim}
lsql-csv -d: 'p=/etc/passwd, 
  "The number of users of "p.7" is "count(p.3)".", by p.7'
\end{verbatim}
The output might be:
\begin{verbatim}
The number of users of /bin/bash is 7.
The number of users of /bin/false is 7.
The number of users of /bin/sh is 1.
The number of users of /bin/sync is 1.
The number of users of /sbin/halt is 1.
The number of users of /sbin/nologin is 46.
The number of users of /sbin/shutdown is 1.
\end{verbatim}

As you can see, string formatting is sometimes very simple with LSQL.

\subsubsection{Arithmetic expression}
So far, we just met all kinds of blocks, and only if block accepts an arithmetic expression and the other accepts a select expression.
What if we needed to run an arithmetic expression inside a select expression? There is a special syntax \textit{\$(...)} for it.

For example:
\begin{verbatim}
lsql-csv -d: '/etc/passwd, $(sin(&1.3)^2 + cos(&1.3)^2)'
\end{verbatim}

It returns something like:
\begin{verbatim}
1.0
1.0
1.0
0.9999999999999999
...
1.0
\end{verbatim}

If we run:
\begin{verbatim}
lsql-csv -d: '/etc/passwd, $(&1.3 >= 1000), sort $(&1.3 >= 1000)'
\end{verbatim}

We get something like:
\begin{verbatim}
false
false
...
false
true
true
...
true
\end{verbatim}

\subsubsection{More complicated join}
Let's see more complicated examples.
\begin{verbatim}
lsql-csv -d: 'p=/etc/passwd g=/etc/group, p.1 g.1, if p.1 in g.4'
\end{verbatim}
This will print all pairs of user and its group excluding the default group. If you know SQL, you can read it as \textit{from /etc/passwd P, /etc/group G select P.1, G.1 where P.1 in G.4}.

How does \textit{in} work? It's one of the basic string level "consist". If A is a substring of B, then \textit{A in B} is true. Otherwise, it is false.

And the output?
\begin{verbatim}
root:root                                                                                                                                                                                                                                                                 
root:wheel                                                                                                                                                                                                                                                                
root:floppy                                                                                                                                                                                                                                                               
root:tape                                                                                                                                                                                                                                                                 
lp:lp                                                                                                                                                                                                                                                                     
halt:root                                                                                                                                                                                                                                                                 
halt:wheel 
\end{verbatim}

\subsubsection{More complicated...}
The previous example doesn't give a very readable output. We can use \textit{group by} to improve it (shortened as \textit{by}).
\begin{verbatim}
lsql-csv -d: 'p=/etc/passwd g=/etc/group, 
  p.1 cat(g.1","), if p.1 in g.4, by p.1'
\end{verbatim}
The output will be something like:
\begin{verbatim}
adm:adm,disk,sys,
bin:bin,daemon,sys,
daemon:adm,bin,daemon,
lp:lp,
mythtv:audio,cdrom,tty,video,
news:news,
\end{verbatim}
It groups all non-default groups of a user to one line and concatenates it delimited by ",".

How can we add default groups too?
\begin{verbatim}
lsql-csv -d: 'p=/etc/passwd g=/etc/group, 
  p.1 cat(g.1","), if p.1 in g.4, by p.1' |
lsql-csv -d: '- /etc/passwd /etc/group, 
  &1.1 &1.2""&3.1, if &1.1 == &2.1 && &2.4 == &3.3'
\end{verbatim}
This will output something like:
\begin{verbatim}
adm:adm,disk,sys,adm
bin:bin,daemon,sys,bin
daemon:adm,bin,daemon,daemon
lp:lp,lp
mythtv:audio,cdrom,tty,video,mythtv
news:news,news
\end{verbatim}

The first part of the command is the same as in the previous example. The second part inner joins the output
of the first part with \textit{/etc/passwd} on the username and \textit{/etc/group} on the default GID number and prints
the output of the first part with an added default group name.



\section{Usage}
Now, if you understand the examples, it is time to move forward to a more abstract description of the language and tool usage.

\subsection{Options}
\begin{verbatim}
-h
--help
\end{verbatim}
Shows short command line help and exits before doing anything else.

\begin{verbatim}
-n
--named
\end{verbatim}
Enables first-line naming convention in csv files. This works only on input files. 
Output is always without first-line column names.

\begin{verbatim}
-dCHAR
--delimiter=CHAR
\end{verbatim}
Changes default primary delimiter. The default value is \textit{,}.

\begin{verbatim}
-sCHAR
--secondary-delimiter=CHAR
\end{verbatim}
Changes default quote char (secondary delimiter). The default value is \textit{"}.

\subsection{Datatypes}
There are 4 datatypes considered: \textit{Bool}, \textit{Int}, \textit{Double}, \textit{String}. 
\textit{Bool} is either true/false, \textit{Int} is at least a 30-bit integer, \textit{Double} double-precision floating point number, and \textit{String} is an ordinary char string.

During CSV data parsing, the following logic of datatype selection is used:
\begin{itemize}
    \item \textit{Bool}, if \textit{true} or \textit{false}
    \item \textit{Int}, if \textit{[0-9]+} matches
    \item \textit{Double}, if \textit{[0-9]+.[0-9]+(e[0-9]+)?} matches
    \item \textit{String}, if none of the above matches
\end{itemize}

\subsection{Joins}
Join means, that you put multiple input files into from block.

Joins always have the time complexity $\mathcal{O}(nm)$. There is no optimization made based on if conditions when you put multiple files into from block.

\subsection{Documentation of language}
\begin{verbatim}
lsql-csv [OPTIONS] COMMAND
      
Description of the grammar
  
  COMMAND -> FROM_BLOCK, REST
                            
  REST -> SELECT_BLOCK, REST 
  REST -> BY_BLOCK, REST   
  REST -> SORT_BLOCK, REST 
  REST -> IF_BLOCK, REST   
  REST ->
  
  FROM_BLOCK -> FROM_SELECTOR FROM_BLOCK
  FROM_SELECTOR ~~> FROM ... FROM //Wildcard and brace expansion
  
  FROM -> FROM_NAME=FROM_FILE OPTIONS
  FROM -> FROM_FILE OPTIONS
  
  OPTIONS -> -dCHAR OPTIONS
  OPTIONS -> --delimiter=CHAR OPTIONS
  OPTIONS -> -sCHAR OPTIONS
  OPTIONS -> --secondary-delimiter=CHAR OPTIONS
  OPTIONS -> -n OPTIONS  
  OPTIONS -> --named OPTIONS
  OPTIONS -> -N OPTIONS   
  OPTIONS -> --not-named OPTIONS
  OPTIONS ->
  
  
  SELECT_BLOCK -> SELECT_EXPR
  BY_BLOCK -> by SELECT_EXPR
  SORT_BLOCK -> sort SELECT_EXPR                                                                                                                              
  IF_BLOCK -> if ARITHMETIC_EXPR
  
                                                                                                                 
  ARITHMETIC_EXPR -> ATOM  
  ARITHMETIC_EXPR -> ONEARG_FUNCTION(ARITHMETIC_EXPR)                                                                                                                                            
  ARITHMETIC_EXPR -> ARITHMETIC_EXPR OPERATOR ARITHMETIC_EXPR
  ARITHMETIC_EXPR -> (ARITHMETIC_EXPR)
  // Logical negation
  ARITHMETIC_EXPR -> ! ARITHMETIC_EXPR
  ARITHMETIC_EXPR -> - ARITHMETIC_EXPR
  
  SELECT_EXPR -> ATOM_SELECTOR SELECT_EXPR                                                                                                                                     
  SELECT_EXPR ->
  
  ATOM_SELECTOR ~~> ATOM ... ATOM //Wildcard and brace expansion
  
  // eg. 1.0, "text", 'text', 1
  ATOM -> CONSTANT
  // eg. &1.1
  ATOM -> COLUMN_NAME
  ATOM -> pi   
  ATOM -> e
  ATOM -> true
  ATOM -> false
  ATOM -> $(ARITHMETIC_EXPR)
  ATOM -> AGGREGATE_FUNCTION(SELECT_EXPR)
  ATOM -> ONEARG_FUNCTION(ARITHMETIC_EXPR)

  // # is not really char:
  // two atoms can be written without space 
  // and will be (string) appended, 
  // if they are separated using quote chars:
  //   left atom must end or right atom must begin with quote char
  // This rule doesn't apply inside ARITHMETIC_EXPR
  ATOM ~~> ATOM#ATOM     

  
  AGGREGATE_FUNCTION -> cat
  AGGREGATE_FUNCTION -> sum
  AGGREGATE_FUNCTION -> count
  AGGREGATE_FUNCTION -> max
  AGGREGATE_FUNCTION -> min
  AGGREGATE_FUNCTION -> avg
  
  //All trigonometric functions in radian
  ONEARG_FUNCTION -> sin
  ONEARG_FUNCTION -> cos
  ONEARG_FUNCTION -> tan
  
  ONEARG_FUNCTION -> asin
  ONEARG_FUNCTION -> acos
  ONEARG_FUNCTION -> atan
  
  ONEARG_FUNCTION -> sinh
  ONEARG_FUNCTION -> cosh
  ONEARG_FUNCTION -> tanh
  
  ONEARG_FUNCTION -> asinh
  ONEARG_FUNCTION -> acosh
  ONEARG_FUNCTION -> atanh
  
  ONEARG_FUNCTION -> exp
  ONEARG_FUNCTION -> sqrt
  
  //The length of the string
  ONEARG_FUNCTION -> size
  ONEARG_FUNCTION -> to_string
  
  ONEARG_FUNCTION -> negate
  ONEARG_FUNCTION -> abs
  ONEARG_FUNCTION -> signum
  
  ONEARG_FUNCTION -> truncate
  ONEARG_FUNCTION -> ceiling
  ONEARG_FUNCTION -> floor
  
  ONEARG_FUNCTION -> even
  ONEARG_FUNCTION -> odd
  
  // A in B means A is a substring of B
  OPERATOR -> in
  
  OPERATOR -> *
  OPERATOR -> **    //general power
  OPERATOR -> ^     //natural power
  OPERATOR -> /
  
  // Integer division truncated towards minus infinity
  // (x div y)*y + (x mod y) == x
  OPERATOR -> div
  OPERATOR -> mod
  
  //Integer division truncated towards 0
  // (x quot y)*y + (x rem y) == x  
  OPERATOR -> quot
  OPERATOR -> rem

  // greatest common divisor
  OPERATOR -> gcd
  // least common multiple
  OPERATOR -> lcm
  
  OPERATOR -> ++    //append
  
  OPERATOR -> +
  OPERATOR -> -
  
  OPERATOR -> <=
  OPERATOR -> >=
  OPERATOR -> <
  OPERATOR -> >
  OPERATOR -> !=
  OPERATOR -> ==
  
  OPERATOR -> ||
  OPERATOR -> &&
\end{verbatim}

Each command is made from blocks separated by a comma. There are these types of blocks.
\begin{itemize}
    \item From block
    \item Select block
    \item If block
    \item By block
    \item Sort block
\end{itemize}

The first block is always from block. If the block after the first block is without a specifier (\textit{if}, \textit{by}, or \textit{sort}), then it is a select block. Otherwise, it is a block specified by the specifier.

From block accept specific grammar (as specified in the grammar description), select, by, and sort block select expression (\textit{SELECT\_EXPR} in the grammar) and if block arithmetic expression (\textit{ARITHMETIC\_EXPR} in the grammar).

Every source file has a number and may have multiple names - assign name, the name given to the source file by \textit{ASSIGN\_NAME=FILE\_PATH} syntax in from block, and  
default name, which is given the path to the file or \textit{-} in case of stdin in from block.

Each column of a source file has a number and may have a name (if the named option is enabled for the given source file).

If the source file with index M (numbering input files from 1) has been given a name XXX, its columns can be addressed by \&M.N or XXX.N, where N is the index of column (numbering columns from 1).  
If the named option is enabled and a column has the name \textit{NAME}, it can also be addressed by \&M.NAME or XXX.NAME.

If there is a collision in naming (two source files have the same name or two columns under the same source file have the same name), then the behavior is undefined.

\subsubsection{Exotic chars}
There are some chars that cannot be in symbol names (column names). For simplicity, we can suppose, they are everything but alphanumerical chars excluding \textit{-}, \textit{.}, \textit{\&} and \textit{\_}.  
Also first char of a symbol name must be non-numerical to not be considered as an exotic char.
Referencing names containing exotic chars without quotes is unsupported.

It is possible to reference columns with names with exotic chars using \textit{`} quote - like \textit{`EXOTIC NAME`}. The source file name is always part of the column name from the syntax perspective of language - it must be inside the quotes.

\subsubsection{Quote chars}
There are 3 quotes (\textit{`}, \textit{"} and \textit{'}) used in Lsql. \textit{"} and \textit{'} are always quoting a string. The \textit{`} quote is used for quoting symbol names.

These chars can be used for fast appending. If two atoms inside \textit{SELECT\_\-EXPR} are written without space and are separated using the quotes, they will be appended. For example, \textit{abc"abc"} means: append column abc to the string abc.

\subsubsection{Constants}
There are 3 types of constants. String, Double, and Int. Everything quoted in \textit{"} or \textit{'} is always String constant. Numbers without \textit{[0-9]+} are considered Int constant and numbers \textit{[0-9]+.[0-9]+} Double constant.

\subsubsection{Operator precedence}
The following list outlines the precedence and associativity of lsql-csv infix operators. The lower the precedence number, the higher the priority.

\begin{enumerate}
    \item \textit{in}, \textit{**}, \textit{\^}
    \item \textit{*}, \textit{/}, \textit{div}, \textit{quot}, \textit{rem}, \textit{mod}, \textit{gcd}, \textit{lcm}
    \item \textit{++}, \textit{+}, \textit{-}
    \item \textit{\textless=}, \textit{\textgreater=}, \textit{\textless}, \textit{\textgreater}, \textit{!=}, \textit{==}
    \item \textit{$\vert\vert$}, \textit{\&\&}
\end{enumerate}

\subsubsection{Select expression}
They are similar to bash expressions\cite{bash-reference-manual}. They are made by atom selector expressions separated by whitespaces. These expressions are expanded, evaluated, and matched to column names, constants, aggregate functions, or arithmetic expressions.

Every atom selector expression can consist
\begin{itemize}
    \item Wildcard (Each wildcard will be expanded to multiple statements during processing)
    \item Bash brace expansion (e.g. {22..25} -\textgreater 22 23 24 25) \cite{bash-reference-manual}
    \item Arithmetic expression in \textit{\$(expr)} format
    \item Quotes \textit{`anything`} to prevent wildcards and expansions
    \item Quotes \textit{"} or \textit{'} to insert string
    \item Call of aggregate function \textit{AGGREGATE\_FUNCTION(next select block)} - there cannot be any space after FUNCTION
    \item Call of single arg function \textit{ONEARG\_FUNCTION(arithmetic expression)} - there cannot be any space after FUNCTION
    \item Constants
    \item Reference to a column name
\end{itemize}

If you want to concatenate strings without \textit{++} operator, you can write: \textit{a.1","a.2}.

Please, keep in mind, that operators must be put inside arithmetic expressions, or they will be matched to a column name or aggregate function.

\subsubsection{Arithmetic expression}
The statement uses mainly classical awk logic.\cite{awk-reference-manual} You can use keywords \textgreater, \textless, \textless=, \textgreater=, ==, $\vert\vert$, \&\&, +, -, *, /,...

\subsubsection{Select blocks}
These blocks determine output. They accept select expressions and are evaluated and printed in a delimitered format.

Every select block must contain at least one reference to the column name, or the behavior is undefined.

Examples of select blocks:
\begin{verbatim}
&1.[3-6]
\end{verbatim}
This will print columns 3, 4, 5, and 6 from the first file.

\begin{verbatim}
ax*.{6..4}
\end{verbatim}
This will print the 6th, 5th, and 4th of all files whose name begins with ax.

\subsubsection{From blocks}
There must be exactly one from block at the beginning of the command. 
The block can contain any files (and \textit{-} specifies standard input). 
You can use any syntax you would otherwise use in bash to select these files (wildcards, expansion,...)\break\cite{bash-reference-manual}. 
You can also overname the file using \textit{NAME=stmt}. If there is more than 1 matching of stmt, the files will be named \textit{(NAME, NAME1, NAME2,...)}.

Example:
\begin{verbatim}
/etc/{passwd,group}
\end{verbatim}

This will select \textit{/etc/passwd} and \textit{/etc/group} files. They can be addressed either as \textit{\&1} or \textit{/etc/passwd}, and \textit{\&2} or \textit{/etc/group}.

If \textit{filename} is put inside \textit{`} quotes, no wildcard or expansion logic will apply to it.

You can also add custom attributes to files in the format \break \textit{FILE -aX --attribute=X -b}. The attributes will be applied to all files which will be matched using \textit{FILE} bash expression.

\paragraph{Possible attributes}

\begin{verbatim}
-n
--named
\end{verbatim}
It means that csv file has the first line with the names of the columns

\begin{verbatim}
-N
--not-named
\end{verbatim}
You can also set the exact opposite. This can be useful if you change the default behavior.

\begin{verbatim}
-dCHAR
--delimiter=CHAR
\end{verbatim}
This changes the primary delimiter.

\begin{verbatim}
-sCHAR
--secondary-delimiter=CHAR
\end{verbatim}
This changes the secondary delimiter char.

Example:
\begin{verbatim}
/etc/passwd -d:
\end{verbatim}

Currently, commas and CHARs, which are also quotes in Lsql, are not supported as delimiters.

\subsubsection{If block}
This block always begins with \textit{if}. They accept arithmetic expressions, which should be convertible to Bool - either String \textit{false}/\textit{true}, Int (\textit{0} \textit{false}, anything else \textit{true}), or Bool. 
Rows with true are printed or aggregated, and rows with false are skipped.

Filtering is done before the aggregation.

You can imagine if block as where clause in SQL.

\subsubsection{By block}
This statement always begins with \textit{by} and the rest of the block is a select expression. There can be only one By block in the whole command.

The by block is used to group the resulting set by the given atoms.

You can imagine by block as the group by clause in SQL. 

There must be at least one aggregate function in the select block if by block is present. Otherwise, behavior is undefined.

If there is an aggregate function present without by block present, aggregation runs over all rows at once.

\subsubsection{Sort block}
This block can be at the end of the command. It begins with the \textit{sort} keyword and the rest is a select expression.

The sort block determines the order of the final output - given atoms are sorted in ascending order.

You can imagine a sort block as the order by clause in SQL.

There can be only one Sort block in the whole command.


\chapter{Developer documentation}
This chapter is for potential developers of the project.

\section{Project building and testing}
The project has two ways of building. One way is through Makefile and the second is through Cabal. 
By running the following command it generates the build folder and \icode{lsql-csv} binary under it.
\begin{verbatim}
make
\end{verbatim}
 
It is necessary to have all Haskell dependencies (\icode{Parsec} ($\geq 3.1$ $<3.2$), \icode{Glob} ($\geq 0.10$ $<0.11$), 
\icode{base} ($\geq 4.9$ $<4.20$), \icode{text} ($\geq 1.2$ $<2.2$) and \icode{containers} ($\geq 0.5$ $<0.8$)) installed.
The package boundaries given are identical to cabal boundaries. Also, it is necessary, that you have \icode{GHC} ($\geq 8 <9.29$) installed.

The second way of building is through cabal, which handles all dependencies for you.

\begin{verbatim}
cabal build
\end{verbatim}


The project unit tests require building through \icode{Makefile} and is called by:
\begin{verbatim}
make test
\end{verbatim}
It should always succeed before any commit is made.

It is possible to generate Haddock developer documentation by calling:
\begin{verbatim}
cabal haddock
\end{verbatim}

The documentation contains comments on all exported functions. The documentation can be alternatively generated
by running:

\begin{verbatim}
make docs
\end{verbatim}
It generates html documentation under the build folder.

The package is also published at \url{https://hackage.haskell.org/package/lsql-csv} in the Hackage public repository. 
The generated documentation is fully browseable there.

\section{\icode{String} vs \icode{Data.Text}}
As there is a long-term discussion in the Haskell community about whether \icode{String} or \icode{Data.Text} should be used as the 
primary representation of text, I would like to emphasize that in this project, \icode{String} is used as a primary
representation for text.

\section{Project layout}

The project is split into:
\begin{enumerate}
  \item a library, which contains almost all the logic and is placed under
       \icode{src} folder of the project

  \item the \icode{main}, which contains one source file
      with \icode{Main}, which is the entry point for \icode{lsql-csv} binary. 
      It parses the arguments, and checks, whether the help optional argument was called or no argument at all was given, 
      and either shows help or further call run from \icode{Lsql.Csv.Main} from the library.
\end{enumerate}

The library is split into 5 different namespaces. Its usage is not strictly defined, but this can be said:
\begin{itemize}
    \item \icode{Lsql.Csv} -- This namespace contains the entry point
    \item \icode{Lsql.Csv.Core} -- This namespace contains the logic of evaluation
    \item \icode{Lsql.Csv.Lang} -- This namespace contains parsers
    \item \icode{Lsql.Csv.Lang.From} -- This namespace contains parsers for the from block
    \item \icode{Lsql.Csv.Utils} -- This namespace contains helper functions
\end{itemize}

\section{Modules}
The following section is a summary of all modules of the library.
\begin{itemize}
    \item \icode{Lsql.Csv.Core.BlockOps} -- This module contains the \icode{Block} definition representing a command block and functions for getting a specific type of blocks from a list of \icode{Block}.
    \item \icode{Lsql.Csv.Core.Evaluator} -- This module contains the evaluator of a \break\icode{lsql-csv} program.
    \item \icode{Lsql.Csv.Core.Functions} -- This module contains the syntactic tree definition and helper functions for its evaluation.
    \item \icode{Lsql.Csv.Core.Symbols} -- This module contains the definition of \icode{Symbol}, \icode{SymbolMap} and helper functions.
    \item \icode{Lsql.Csv.Core.Tables} -- This module contains the definition of \icode{Value}, \icode{Table}, and \icode{Column}, classes over them, and functions for manipulation of them.
    \item \icode{Lsql.Csv.Lang.Args} -- A module for command line argument parsing.
    \item \icode{Lsql.Csv.Lang.BlockChain} -- This module contains the main parser of blo\-cks other than the from block.
    \item \icode{Lsql.Csv.Lang.BlockSeparator} -- This module contains the preprocessor parser, which splits a command into a list of strings - one string per one block.
    \item \icode{Lsql.Csv.Lang.Options} -- This module implements the common \icode{Option} type for from blocks and command-line optional arguments representation, and its parsers.
    \item \icode{Lsql.Csv.Lang.Selector} -- This module implements the selector expression parser and the arithmetic expression parser.
    \item \icode{Lsql.Csv.Lang.From.Block} -- This module contains the from block parser. It loads the initial \icode{SymbolMap}.
    \item \icode{Lsql.Csv.Lang.From.CsvParser} -- This module contains the \icode{CsvParser} called by the \icode{parseFile}, which loads input CSV files.
    \item \icode{Lsql.Csv.Main} -- This module contains the starting point for a \icode{lsql-csv} evaluation.
    \item \icode{Lsql.Csv.Utils.BracketExpansion} -- This module contains the curly bra\-cket (braces) expansion implementation.
    \item \icode{Lsql.Csv.Utils.CsvGenerator} -- This module contains the CSV generator for the output.
\end{itemize}

\section{Entry point}
The entry point is in module \icode{Lsql.Csv.Main} in function \icode{run}. 
The function first calls preprocessor in \icode{Lsql.Csv.Lang.BlockSeparator}, which splits command into a list of strings. Then \icode{SymbolMap} with input is loaded using \icode{Lsql.\allowbreak Csv.Lang.From.Block} and after that
rest of the blocks are parsed using \icode{Lsql.\allowbreak Csv.\allowbreak Lang.\allowbreak BlockChain}. The program is then evaluated using
\icode{Lsql.\allowbreak Csv.Core.\allowbreak Evaluator} and finally, the output generated by \icode{Lsql.\allowbreak Csv.Utils.CsvGenerator}.




\chapter{Analysis}
Why the language have been made the way it is? Why it is so inspired by SQL and is not just next implementation of SQL?
Why it is implemented the way it is?
Further chapter is about design decision of the language.


\section{Why SQL at first place?}
Why we talk so much about SQL at first place? Since it was introduced in 1970s \cite{enwiki-sql}, it has
become de facto standard for many major databases. Just for illustration name a few of them.
\begin{itemize}
    \item Oracle DB have used it since 1979 as a first commercially available implementation \cite{oracle-db-sql}. 
    \item MySQL have used it since 1994, since its original development started \cite{enwiki-mysql} 
    \item PostgreSQL have used it since 1996, since it was created \cite{postgres-birthday}. 
    \item MSSQL have used it since 1989, since its initial release \cite{enwiki-mssql}. 
\end{itemize}

SQL is so much known, that there is a widely used term NoSQL databases as databases opposed to SQL databases.

The main point of making language inspired by SQL is that it brings the advantage of getting large user base which only needs to understand the difference between SQL and the new language to start using the new language. This is the starting point, from which we further argument about design decisions made.


\section{Why not implement just another SQL for CSV files?}
As mentioned in the introduction, standard SQL is a typed language (Every data column and data value belongs to some data type) \cite{ISO9075-2023}, 
which implied many design choices. Furthermore, SQL itself require large amount of text to be written,
before it can be executed. 

One of the ambition of \textit{lsql-csv} is to allow a user to write shorter queries to get the result. For the example,
consider

\begin{verbatim}
select A.dataX from data A where A.dataX > 1000
\end{verbatim}

This simple SQL query shows $dataX > 1000$ from table \textit{data}. 
Now, if we have a CSV file $data.txt$, where we know, that \textit{dataX} is a second column, the same query can 
be written with \textit{lsql-csv} as

\begin{verbatim}
data.txt, &1.2, if &1.2 > 1000
\end{verbatim}

The length difference is about $35\%$ off the original query. It is simply said the another reason, 
why we will not just implement another SQL implementation.

\section{Why number references?}
Where the $35\%$ difference happened. One and the main reason, why is, that we allowed referencing $dataX$ as \textit{.2}.
This is also possible due to nature of CSV file, where columns have their index
\footnote{Do not confuse with SQL database index for performance}. 
Normally, in SQL database, indexes of columns are not considered as something, which should decide about query meaning.
This is because the SQL database itself may change, new columns may be added or removed. Just because somebody removed a column, 
it is not wanted, that developers needed to change some of the queries so they comply with the new database layout.

On the other hand, CSV files usually are not altered as much as databases are. 
If developer (or an user) is not sure about the columns of the CSV file, it is one of the first signs, that
he or she should use rather SQL database than simple CSV file.

Also \textit{lsql-csv} is rather a tool for daily life and simple scripts rather than tool for development of medium sized or
large sized project, like SQL is. This adds much more flexibility in what can language do (like number references).

\section{Why renaming standard keywords?}
Why have we renamed \textit{where} for \textit{if}, \textit{group by} for \textit{by} and \textit{sort by} for \textit{sort}? Simply said, because the renamed variants are shorter and still do not block user in understanding, what the query does.

\section{Why some features like descending sort are missing?}
\textit{lsql-csv} is a shell utility and as such, it tries to comply with UNIX philosophy. 
The summarized version from Doug McIlroy is: `This is the Unix philosophy: Write programs that do one thing and do it well. Write programs to work together. Write programs to handle text streams, because that is a universal interface.' \cite{unix-philosophy}.

The tool for reverting the order of sorted output already exist: \textit{tac}. Piped together it creates the wanted output.

Similar case, the function for second filtering of grouped by (in SQL named \textit{having}) 
output have been not added, because the wanted output
may be received by piping the output with another instance of \textit{lsql-csv}.

And why there is no support for creating output with first line names of columns? Because the wanted output
may be simply made with usage of \textit{echo} called before \textit{lsql-csv} if needed.

\section{Why blocks are delimited by comma?}
The author thinks, it is more readable like this. Developers in SQL usually use upper case and new line writing to delimiter the blocks and comma is just more easy to write than switching on and off caps lock, holding shift key or making multiline input.

\section{Why there are two types of expression?}
Why there are arithmetic and select expression with different grammars?
We think, that addition of special select expression further allows the user to write shorter queries.
It allows us to introduce wildcard and curly brackets expansion for the user, which will not be possible otherwise.

\section{Why there is support only for cross join and not other types of join?}
The tool is supposed to be simple. We recommend user to import CSV data to SQL database and use standard SQL if he needs more complicated joins.

The other reason is, that CSV have no standardized NULL value, which is needed for left or right outer joins.

\section{Why there is no package used for CSV parsing and generating?}
Very simply said, to limit number of dependencies. The more dependencies is used, the harder is to compile it, maintain it and add to any Linux distribution.

As CSV parsing is not hard job, it was decided so.

\section{Why \textit{String} is used as primary text representation?}
As the performance gain from using \textit{Data.Text} would not be significant and \textit{String} does not add any more complexity to the code like \textit{Data.Text} does, 
it was decided that \textit{String} will be primary text representation.

\section{Why joins have $\mathcal{O}(nm)$ complexity?}
\textit{lsql-csv} is a tool for small dataset data querying. As such, it implements only simple algorithm for joining the tables -- cross join.

For larger dataset users are encouraged to use standard SQL database.





\chapter{Alternative Solutions}
What other solutions are there for the given problem? What other approaches can we use, when we are dealing with queries over CSV files? 
The following chapter is about alternative approaches to the problem.

\section{Using SQL database}
The first obvious solution is importing the dataset to some standard SQL database and doing the queries over it. 
This approach requires the definition of the schema into which the data will be imported. This is an overhead, which isn't
always advantageous to pay as we might need only a simple query to be run over it.
The data import might take also a much longer time than a simple \icode{lsql-csv} invocation.

But it might be a very favorable solution, if we need large dataset joins, need a complex query execution or a large amount of simple queries run over it.

By using a standard SQL database you gain the advantage of better performance, typed datasets, indexes, and a larger number of built-in functions.

\section{Using standard Unix tools}
It is possible to do a large amount of work just by using \icode{awk}, \icode{join}, \icode{sort}\ldots{}
For example:

\begin{verbatim}
lsql-csv -d: '-, &1.*, if &1.3 >= 1000' </etc/passwd
\end{verbatim}

This query might be rewritten to \icode{awk}:
\begin{verbatim}
awk -F: '{ if($3 >= 1000){ print $0 }}' </etc/passwd
\end{verbatim}

The main advantage of \icode{lsql-csv} is that it handles some more complex queries more easily. For example:
\begin{verbatim}
lsql-csv -d: '/etc/{passwd,group}, &1.1 &2.1, if &1.4 == &2.3'
\end{verbatim}

This is a simple join query. When written using standard Unix tools, it is:
\begin{verbatim}
sort -t: -k3,3 /etc/group >/tmp/group.sort
sort -t: -k4,4 /etc/passwd >/tmp/passwd.sort
join -t: -14 -23 /tmp/passwd.sort /tmp/group.sort | cut -d: -f2,8 
\end{verbatim}

As demonstrated, the \icode{lsql-csv} variant is 
more readable and shorter\footnote{It is possible, that it can be written in shorter and more readable form, but probably not more than the \icode{lsql-csv} variant.}.

It should be also noted, that an \icode{lsql-csv} join has $\mathcal{O}(nm)$ time complexity, while standard Unix tools have for join written above
$\mathcal{O}(n\log n + m\log m)$ time complexity, so for larger datasets, it might be more beneficial to use them.

\section{By using SQL implementation for CSV files}
There are many projects implementing SQL on CSV files. For example:
\begin{itemize}
    \item \icode{q} \cite{q}
    \item \icode{CSV SQL} \cite{csv-sql} 
    \item \icode{trdsql} \cite{trdsql}
    \item \icode{csvq} \cite{csvq}
\end{itemize}

It is possible to use them to do the job directly with SQL. 

The advantage of it is you do not have to learn a new language if you already know standard SQL. 
The disadvantage is that the queries will be probably longer than would be with \icode{lsql-csv}.

\section{By using general-purpose programming language}
It is not hard to parse and process CSV files with a general-purpose programming language (for example Python).

The advantage of this solution is a much greater flexibility of what you can do with the CSV files.
The large disadvantage is, that it will take too much code to be written for any query.



\chapter*{Conclusion}
\addcontentsline{toc}{chapter}{Conclusion}

The goal of this thesis was to create a new tool for small CSV file data querying from a shell with short queries.
It ought to be simple, comply with Unix philosophy, and use more lapidary language than SQL.
It ought to enable its users to quickly write simple queries directly to the terminal.

Our new tool \icode{lsql-csv} implements a new language LSQL similar to
SQL, specifically designed for working with CSV files in a shell. It makes it possible to work with small CSV files like with a read-only relational database.

In some of the use cases, the tool provides shorter and more readable commands then you would get by using standard Unix tools and the 
language is more lapidary than SQL.

The \icode{lsql-csv} package has been uploaded to the Hackage public package repository and therefore is easily installable. 
It provides automatic Haddock documentation, which is browseable from there. 
It has also been published on GitHub with the \icode{README.md} file containing the full user documentation including the tutorial and grammar description.



%%% Bibliography
%%% Bibliography (literature used as a source)
%%%
%%% We employ bibTeX to construct the bibliography. It processes
%%% citations in the text (e.g., the \cite{...} macro) and looks up
%%% relevant entries in the bibliography.bib file.
%%%
%%% The \bibliographystyle command selects, which style will be used
%%% for references from the text. The argument in curly brackets is
%%% the name of the corresponding style file (*.bst). Both styles
%%% mentioned in this template are included in LaTeX distributions.

% \bibliographystyle{plainnat}    %% Author (year)
\bibliographystyle{unsrt}     %% [number]

\renewcommand{\bibname}{Bibliography}

%%% Generate the bibliography. Beware that if you cited no works,
%%% the empty list will be omitted completely.

\bibliography{bibliography}

%%% If case you prefer to write the bibliography manually (without bibTeX),
%%% you can use the following. Please follow the ISO 690 standard and
%%% citation conventions of your field of research.

% \begin{thebibliography}{99}
%
% \bibitem{lamport94}
%   {\sc Lamport,} Leslie.
%   \emph{\LaTeX: A Document Preparation System}.
%   2nd edition.
%   Massachusetts: Addison Wesley, 1994.
%   ISBN 0-201-52983-1.
%
% \end{thebibliography}



%%% Abbreviations used in the thesis, if any, including their explanation
%%% In mathematical theses, it could be better to move the list of abbreviations to the beginning of the thesis.

\openright
\end{document}
